\chapter{System Architecture}
\label{ch:architecture}

\section{The system, end to end}

\begin{figure}[h]
\centering
\begin{tikzpicture}[
  box/.style={draw=olkSlateMuted, thick, rounded corners=2pt, minimum height=1.1cm, align=center, font=\sffamily\small},
  svc/.style={box, fill=olkTeal!8},
  ext/.style={box, fill=olkSlate!5, draw=olkGray},
  lbl/.style={font=\sffamily\scriptsize\color{olkGray}},
  arr/.style={-Stealth, thick, color=olkSlateMuted}
]
  \node[box, fill=olkSlateLight, text=white, minimum width=3.2cm] (browser) {Browser\\[1pt]\footnotesize (reader in Kashmir)};

  \node[svc, minimum width=6.8cm, minimum height=3.6cm, below=1.4cm of browser, yshift=-0.3cm] (vercel) {};
  \node[lbl, above right] at (vercel.north west) {\textbf{Vercel} — Next.js App Router};

  \node[box, fill=white, minimum width=2.9cm, minimum height=0.9cm] (rsc) at ($(vercel.center)+(-1.6cm,0.9cm)$) {Server Components\\[-1pt]\footnotesize \& Server Actions};
  \node[box, fill=white, minimum width=2.6cm, minimum height=0.9cm] (proxy) at ($(vercel.center)+(1.9cm,0.9cm)$) {Proxy\\[-1pt]\footnotesize (\texttt{proxy.ts})};
  \node[box, fill=white, minimum width=2.9cm, minimum height=0.9cm] (cc) at ($(vercel.center)+(-1.6cm,-0.7cm)$) {Client Components\\[-1pt]\footnotesize (browser JS bundle)};
  \node[box, fill=white, minimum width=2.6cm, minimum height=0.9cm] (api) at ($(vercel.center)+(1.9cm,-0.7cm)$) {Route Handlers\\[-1pt]\footnotesize \pathname{/api/digest}, \pathname{/api/books}};

  \node[svc, minimum width=6.6cm, minimum height=3.3cm, below=1.6cm of vercel] (supabase) {};
  \node[lbl, above right] at (supabase.north west) {\textbf{Supabase} (local: Docker; later: self-hosted in Kashmir)};

  \node[box, fill=white, minimum width=2.7cm] (pg) at ($(supabase.center)+(-1.5cm,0.75cm)$) {Postgres\\[-1pt]\footnotesize + Row-Level Security};
  \node[box, fill=white, minimum width=2.3cm] (auth) at ($(supabase.center)+(1.6cm,0.75cm)$) {GoTrue Auth};
  \node[box, fill=white, minimum width=2.7cm] (storage) at ($(supabase.center)+(-1.5cm,-0.75cm)$) {Storage\\[-1pt]\footnotesize (\texttt{book-covers})};
  \node[box, fill=white, minimum width=2.3cm] (rt) at ($(supabase.center)+(1.6cm,-0.75cm)$) {Realtime};

  \node[ext, right=2.3cm of vercel, minimum width=2.6cm] (resend) {Resend\\[-1pt]\footnotesize transactional email};
  \node[ext, below=0.5cm of resend, minimum width=2.6cm] (olib) {Open Library API\\[-1pt]\footnotesize (ISBN lookups)};

  \draw[arr] (browser) -- (vercel.north);
  \draw[arr] (vercel) -- (supabase);
  \draw[arr] (api.east) -- ++(0.7,0) |- (resend.west);
  \draw[arr] (cc.east) -| ($(olib.west)+(-0.4,0)$) -- (olib.west);
\end{tikzpicture}
\caption{OLK's runtime shape. The Client-Component box is the only part of Vercel's box that actually executes in the reader's browser; everything else in that box runs on Vercel's servers. The Supabase box is a single Docker Compose stack in local development and, today, Supabase's hosted cloud in production — the long-term goal (Chapter~\ref{ch:welcome}) is for that same box to be a physical server in Kashmir.}
\label{fig:architecture}
\end{figure}

\section{Server Components, Client Components, and where the boundary actually falls}

Every file under \pathname{src/app} is a Server Component by default. A file becomes a Client Component only if its first line is the string \code{"use client"}. This is not a stylistic choice per file — it changes what APIs are available and where the code physically runs:

\begin{itemize}
  \item \textbf{Server Components} (no directive) render on Vercel's servers, can call \code{await createClient()} from \pathname{src/lib/supabase/server.ts} directly, and never ship their data-fetching code to the browser. \pathname{src/app/page.tsx} (the landing page) and every \pathname{layout.tsx} in this project are Server Components.
  \item \textbf{Client Components} (\code{"use client"}) render once on the server for the initial HTML, then hydrate and re-render in the browser. They use \pathname{src/lib/supabase/client.ts}'s \code{createClient()}, which talks to Supabase directly over HTTPS from the reader's browser — meaning every query a Client Component makes is subject to Row-Level Security as the logged-in user, never as an elevated role. Nearly every interactive page in the dashboard (\pathname{browse}, \pathname{my-books}, \pathname{requests}, \pathname{messages/[id]}, all of \pathname{admin/*}) is a Client Component, because they need \code{useState}/\code{useEffect} for interactivity.
\end{itemize}

\begin{olknote}
One deliberate exception to "Client Components talk to Supabase directly": \pathname{browse/page.tsx}'s plain, unfiltered book list goes through \pathname{/api/books} (the Route Handler in Figure~\ref{fig:architecture}) instead, specifically so its result can be cached and shared across every visitor for a few seconds rather than re-queried per request (Chapter~\ref{ch:proxy}). Any search or filter still bypasses that route and queries Supabase directly, since a filtered result isn't the same for every visitor and so isn't safe to cache under a shared key. This is the general rule for adding a shared cache anywhere else in this codebase: only a query whose result is identical no matter who's asking is a candidate.
\end{olknote}

\begin{olknote}
This project's most complex pages are almost entirely Client Components, even though a "fetch on the server, render once" Server Component would often be more efficient. That is a real, current characteristic of the codebase, not a recommendation — see Chapter~\ref{ch:pages} for page-by-page detail, and treat "should this be a Server Component instead" as a legitimate refactor to propose, not a rule that's already been decided against.
\end{olknote}

\section{Where the Proxy sits}

Every request to the app — before any Server Component runs, before any HTML is generated — passes through \pathname{src/proxy.ts}. It refreshes the Supabase auth session cookies (required because \code{@supabase/ssr} needs a chance to rotate tokens on every request), checks \code{platform\_settings} feature flags and maintenance mode, redirects unauthenticated visitors away from a fixed list of protected route prefixes, and redirects banned users or non-admins away from routes they should not reach. Chapter~\ref{ch:proxy} covers this file line by line.

\section{Local development vs. production, side by side}

\begin{center}
\begin{tabularx}{\textwidth}{@{}l Y Y@{}}
\toprule
& \textbf{Local development} & \textbf{Production (today)} \\
\midrule
Frontend & \code{next dev} on \code{localhost:3000} & Vercel build of the same repository \\
Database \& Auth & Docker containers started by \code{supabase start}, at \code{127.0.0.1:54321}/\code{:54322} & Supabase's hosted cloud project \\
Config source & \pathname{.env.local} pointed at \code{127.0.0.1} & Environment variables set in the Vercel dashboard \\
Email & Silently disabled (no \code{RESEND\_API\_KEY} set) & Live, via Resend \\
Migrations & \code{supabase db reset} replays every file in \pathname{supabase/migrations} & Applied once, in order, via \code{supabase db push} or the linked project \\
\bottomrule
\end{tabularx}
\end{center}

The two columns are deliberately structured to be as close to identical as possible — the same migrations, the same RLS policies, the same functions — so that "it works locally" is a meaningful, trustworthy signal before you ever touch the hosted database. Chapter~\ref{ch:local-env} is the practical walkthrough of the left column.

\begin{olktask}
Pick one concrete action — requesting a book from Browse is a good one — and trace it through Figure~\ref{fig:architecture} by name, using only this chapter and the cross-references it points to, \emph{before} reading the chapters that cover each hop in detail. Which box first receives the click (Client Components — \pathname{browse/page.tsx} is \code{"use client"}, Chapter~\ref{ch:pages})? What talks to Postgres, and as which role (Chapter~\ref{ch:rls})? What inside the Supabase box might reject the request even though your code is correct (Chapter~\ref{ch:functions})? Does the Proxy box get involved for this specific action, or only for the page load that preceded it (Chapter~\ref{ch:proxy})? Write the chain of boxes down, left to right, then go verify it as you read those chapters. Getting this wrong on the first attempt and correcting it teaches you the system's shape far faster than reading the right answer first would.
\end{olktask}
